\documentclass{article}
\usepackage[preprint]{neurips_2025}
\usepackage[T1]{fontenc}
\usepackage{amsmath,amssymb,booktabs,tabularx}
\usepackage{microtype}
\usepackage[hidelinks]{hyperref}
\usepackage{url}
\hypersetup{pdftitle={Asymmetric Process-Reward Adaptation for Mathematical Reasoning},
  pdfauthor={Hongyi Ding; Jasmine Jiang; Yang Zeng}, pdfsubject={ECE285 Project Proposal}}
\title{Asymmetric Process-Reward Adaptation\\for Mathematical Reasoning}
\author{
Hongyi Ding \\
\And
Jasmine Jiang \\
\And
Yang Zeng \\
}
\newcommand{\E}{\mathbb{E}}
\newcommand{\sg}{\operatorname{sg}}
\newcolumntype{L}{>{\raggedright\arraybackslash}X}
\begin{document}
\maketitle
\begin{abstract}
Correct answers can conceal faulty derivations. We propose asymmetric process reward model (PRM) adaptation: retain failure supervision and public step labels while omitting positive outcome supervision. Four matched GRPO experiments test whether this improves mathematical reasoning.
\end{abstract}

\section{Motivation and research question}
Mathematical RL needs reliable rewards to favor valid reasoning over lucky answers. Yet answer-correct solutions can contain errors \citep{zheng2025}, so outcome labels may certify flawed derivations. We ask: \textbf{can withholding positive outcome supervision improve downstream policy accuracy?} Reduced false certification must outweigh lost supervision. A useful rule would adapt public PRMs without new annotations; diagnostics characterize failures without uniquely attributing them to label noise or policy drift.

\section{Related work and contribution}
PRM800K provides costly human step labels on previously sampled solutions \citep{lightman2023}; their coverage need not track a changing policy. PRIME updates implicit rewards online from outcome labels \citep{cui2025}, which do not identify valid individual steps. PURE's min-form credit assignment mitigates but does not eliminate PRM reward hacking \citep{cheng2025}. LCA uses softmax-weighted pooling and tests controlled answer-correct/process-wrong label noise \citep[App.~E.3]{jia2026}. Our contribution is \textbf{withholding weak positive labels while retaining step supervision}, tested through downstream GRPO. C adopts LCA-inspired pooling, not its full method; D--C isolates the removed loss term.

\section{Proposed method}
\textbf{Offline adaptation.} Sample four solutions on each of 512 training problems from the initial policy; retain complete, unambiguously verified trajectories. For native step-boundary correctness scores $p_{\phi,t}$, set $e_{\phi,t}=1-p_{\phi,t}$. Following LCA,
\begin{equation}
 w_t=\frac{\exp(e_{\phi,t}/T)}{\sum_j\exp(e_{\phi,j}/T)},\qquad
 E_\phi(\tau)=\sum_t\sg(w_t)e_{\phi,t},\qquad T=0.2 .
\end{equation}
Here $\sg$ stops gradients through weights. Define $L_{\rm fail}=\E_{\tau^-}[-\log E_\phi(\tau^-)]$ and $L_{\rm success}=\E_{\tau^+}[-\log(1-E_\phi(\tau^+))]$, taking separate means over failures/successes. A shared set of 2,048 public positive/negative step labels supplies mean binary cross-entropy $L_{\rm step}$:
\begin{equation}
 L_c=L_{\rm fail}+cL_{\rm success}+L_{\rm step},\qquad
 c_C=1,\quad \boxed{c_D=0\ \text{(ours)}} .
\end{equation}
Failure labels weakly indicate an error without locating it. Step labels discourage blanket rejection but cannot guarantee it is avoided.

\newpage
\noindent\textbf{Policy training.} Freeze each PRM throughout RL. For complete, parseable responses in B--D, compute
\begin{equation}
 R_i=y_i+0.2\min_t p_{\phi,t}(\tau_i),\qquad
 \hat A_i=\frac{R_i-\bar R}{\sigma_R+10^{-4}} .
\end{equation}
Here $y_i$ is binary correctness; $\bar R,\sigma_R$ denote the group mean/sample standard deviation. A uses $R_i=y_i$; invalid/truncated outputs get zero in all arms. GRPO applies detached advantages to completion tokens \citep{shao2024}. For valid same-outcome groups, \mbox{$\hat A_i=(m_i-\bar m)/(\sigma_m+5\!\times\!10^{-4})$}, where $m_i=\min_t p_{\phi,t}$: 0.2 does \textbf{not} bound normalized PRM influence. Unlike A, B--D can reward complete wrong answers over invalid outputs; C/D share this rule.

\section{Experimental design and expected findings}
\begin{center}
\small
\renewcommand{\arraystretch}{1.12}
\begin{tabularx}{\linewidth}{@{}l L L@{}}
\toprule
\textbf{Arm} & \textbf{Reward and PRM preparation} & \textbf{Question tested} \\
\midrule
A & Outcome only; no PRM & Reference GRPO performance \\
B & Outcome + original frozen PRM & B--A: value of an existing PRM \\
C & Outcome + PRM adapted with $c=1$ & C--B: complete adaptation package \\
\textbf{D (ours)} & Outcome + PRM adapted with $c=0$ & \textbf{D--C: value of the proposed change} \\
\bottomrule
\end{tabularx}
\end{center}
\textbf{Models and controls.} Use Qwen2.5-Math-1.5B (policy LoRA rank 16), the H4 1.5B PRM (PRM LoRA rank 8) \citep{h4prm}, and TRL 0.23.1/PEFT GRPO with group scaling \citep{trl}. Outcome-only execution is verified on one 80\,GB A100. A \textbf{PRM-only control S} uses $L_{\rm step}$ with matched initialization, anchor batches and updates: S--B tests anchor adaptation; D--S tests failure supervision. No RL arm is added.

\textbf{Pre-RL check.} Score shared initial-policy rollouts on 64 validation questions (eight each) with B/C/D/S; audit missing updates. Before B--D RL, fix replication using C--D advantages: if sign agreement is at least 99\% and every group's RMS difference is at most 0.05, use seed 17 only across A--D; otherwise use 17/42/83. This is a budget heuristic, not evidence of equivalence; test accuracy and ProcessBench never choose replication.

\textbf{Data.} Use MATH (6,993 train/500 validation) \citep{hendrycks2021}, MATH500 and GSM8K tests \citep{cobbe2021}. Both tests have zero exact normalized overlap with the public PRM800K inventory. Full MATH overlaps PRM800K \citep{prmdata} and is supplementary. Exclude validation/test/diagnostic questions from new anchors; unknown inherited exposure and paraphrases remain limitations.

\textbf{Budget and matching.} C/D share PRM initialization, batches, optimizer and updates; only $c$ changes, without loss re-averaging. All arms share policy initialization, prompt order and decoding: eight rollouts, temperature 1.0, 2,048-token cap. Fix \textbf{24M completion tokens/run}, with 6M/12M/24M milestones: 96M or 288M total under the gate. Count failures/truncations; report updates, prompts, lengths and GPU hours including PRM/precheck costs.

\textbf{Evaluation and power.} Use identical greedy decoding without reranking. Primary: \textbf{D--C MATH500 accuracy at 24M}; D--A/B tests practical value. Report paired-seed differences and 95\% question-cluster bootstrap intervals retaining all seeds together; mean/std only with three seeds. Single-seed results are exploratory. At an assumed 10\% disagreement rate, the approximate single-seed 80\%-power detectable difference is 4 percentage points (two-sided $\alpha=.05$). Small effects may remain unresolved; milestones do not select the endpoint.

\textbf{Mechanism and risks.} ProcessBench MATH diagnostics include first-error accuracy, correct-chain acceptance, their harmonic mean at threshold 0.5, and valid-chain AUROC among answer-correct solutions (406 valid/94 flawed). Report paired question-bootstrap intervals for metrics/differences and overlap strata; small uncertain differences remain inconclusive. On shared rollouts, report D--C rewards/advantages within fixed all-correct/all-wrong/mixed groups. Online group rates are descriptive, not causal subgroup effects. Flag invalid outputs; monitor steps and missing delimiters with a fixed double-newline splitter. Diagnostics never tune PRMs or establish policy gains alone.

\textbf{Tentative timeline.} Week 1: buffer, anchors, PRM controls and precheck. Weeks 2--3: C/D first, then remaining A/B for seed 17; add paired seeds only as predeclared. Week 4: diagnostics, statistics and report; disclose unfinished runs.
\clearpage
\begin{thebibliography}{11}
\bibitem[Shao et al.(2024)]{shao2024}
Zhihong Shao et al. DeepSeekMath: Pushing the Limits of Mathematical Reasoning in Open Language Models. 2024.
\url{https://arxiv.org/abs/2402.03300v3}.

\bibitem[Zheng et al.(2025)]{zheng2025}
Chujie Zheng et al. ProcessBench: Identifying Process Errors in Mathematical Reasoning. ACL, 2025.
\url{https://arxiv.org/abs/2412.06559v4}.

\bibitem[Lightman et al.(2023)]{lightman2023}
Hunter Lightman et al. Let's Verify Step by Step. 2023.
\url{https://arxiv.org/abs/2305.20050}.

\bibitem[Cui et al.(2025)]{cui2025}
Ganqu Cui et al. Process Reinforcement through Implicit Rewards. 2025.
\url{https://arxiv.org/abs/2502.01456v2}.

\bibitem[Cheng et al.(2025)]{cheng2025}
Jie Cheng et al. Stop Summation: Min-Form Credit Assignment Is All Process Reward Model Needs for Reasoning. NeurIPS, 2025.
\url{https://arxiv.org/abs/2504.15275v3}.

\bibitem[Jia et al.(2026)]{jia2026}
Tianyu Jia et al. The Weakest Link Tells It All: Outcome-Supervised Process Reward Modeling via Learnable Credit Assignment. 2026.
\url{https://arxiv.org/abs/2606.27739v2}.

\bibitem[Hendrycks et al.(2021)]{hendrycks2021}
Dan Hendrycks et al. Measuring Mathematical Problem Solving With the MATH Dataset. NeurIPS, 2021.
\url{https://arxiv.org/abs/2103.03874v2}.

\bibitem[Cobbe et al.(2021)]{cobbe2021}
Karl Cobbe et al. Training Verifiers to Solve Math Word Problems. 2021.
\url{https://arxiv.org/abs/2110.14168v2}.

\bibitem[Hugging Face H4(2025)]{h4prm}
Hugging Face H4. Qwen2.5-Math-1.5B-Instruct-PRM-0.2. 2025.
\href{https://huggingface.co/HuggingFaceH4/Qwen2.5-Math-1.5B-Instruct-PRM-0.2}{Official model card and scoring interface}. Accessed Oct.\ 20, 2026.

\bibitem[Hugging Face(2026)]{trl}
Hugging Face. TRL 0.23.1: \href{https://huggingface.co/docs/trl/v0.23.1/grpo_trainer}{GRPO Trainer documentation}. Accessed Oct.\ 20, 2026.

\bibitem[OpenAI(2023)]{prmdata}
OpenAI. PRM800K dataset repository: MATH splits and answer grading. 2023. Accessed Oct.\ 20, 2026.
\url{https://github.com/openai/prm800k}.
\end{thebibliography}
\end{document}
